% back_conclusion
\clearpage
\phantomsection
\section*{\bgen{ЗАКЛЮЧЕНИЕ}{CONCLUSION}}
\addcontentsline{toc}{section}{\bgen{ЗАКЛЮЧЕНИЕ}{CONCLUSION}}

Целта, формулирана в раздел~\ref{sec:goal}, съдържа две части: да се разработи архитектура, при
която броят на слушащите дескриптори не зависи от броя на обслужваните протоколи, и да се
установи цената на това решение чрез измерване. Заключението разделя двете, защото те са в
различно състояние.

Първата част е изпълнена. Втората е подготвена и не е проведена, и причината за това не е
пропуск, а свойство на наличната машина. Разделът по-долу казва точно кое от какъв вид е.

\subsection*{\bgen{По задачите}{On the tasks}}

\begin{table}[H]
\centering
\caption{Състояние на задачите от раздел~\ref{sec:goal}}
\label{tab:task-status}
\begin{tabular}{@{}L{0.6cm}L{7.6cm}L{5.8cm}@{}}
\toprule
\textbf{№} & \textbf{Задача} & \textbf{Състояние} \\
\midrule
1 & Формална постановка на минимизацията & Доказана, раздел~\ref{sec:formal} \\
2 & Разпознаване на протокола след приемане & Реализирано, раздел~\ref{sec:demux} \\
3 & Преносимо разпределяне на QUIC връзки & Реализирано, раздел~\ref{sec:cid} \\
4 & Отложено получаване с типизиран договор & Реализирано, раздел~\ref{sec:deferred} \\
5 & Втори механизъм за вход и изход под Linux & Реализирано, раздел~\ref{sec:iobackends} \\
6 & Методика за измерване & Формулирана, глава~V \\
7 & Провеждане на измерването & \textbf{Не е проведено} \\
\bottomrule
\end{tabular}
\end{table}

Шест от седемте задачи са изпълнени и всяка от тях сочи мястото, на което може да бъде
проверена. Седмата не е, и това е единственото място в работата, където разликата между
„направено“ и „обосновано“ е съществена.

\subsection*{\bgen{По хипотезите}{On the hypotheses}}

Нито една от четирите хипотези не е потвърдена и нито една не е отхвърлена.

Това е твърдение за машината, а не за хипотезите. Правилата за допустимост от глава~V
отказват запис на резултат при открита виртуализация, и този отказ е сработил: наличната
среда се разпознава като виртуализирана и следователно не може да произведе цитируемо число.
Механизмът е част от методиката и е предвиден да сработва точно в такъв случай.

Алтернативата би била да се съобщят числа от нея с уговорка в бележка под линия. Тя се
отхвърля, защото уговорката не пътува заедно с числото: цитира се числото. Затова
глава~VI съдържа маркери на липсващи стойности, а не правдоподобни стойности.

Всяка от четирите хипотези остава формулирана така, че да може да бъде отхвърлена, и
измерването, което би я отхвърлило, е записано заедно с нея. Кампанията е дефинирана
изцяло: остава да бъде изпълнена върху невиртуализирана машина под Linux.

\subsection*{\bgen{Какво е установено независимо от измерването}{What is established independently of measurement}}

Част от твърденията в работата не зависят от кампанията, защото не са твърдения за
производителност. Те се изброяват тук, за да не бъдат подминати заедно с онези, които зависят.

\paragraph{Предложението за дескрипторите е доказано, а не измерено.} Равенството
\(L = W \cdot |T|\) следва от несечащите се множества от начални октети и от предположенията,
изброени в раздел~\ref{sec:formal}. Броенето на дескриптори на работещ процес го потвърждава
като наблюдение, но не е негово основание.

\paragraph{Трите протокола се обслужват от един екземпляр и един номер на порт.} Проверено е
от край до край, включително установяване на връзка по HTTP/3 срещу еталонен клиент.
Твърдението е за функция, не за скорост, и не се променя от резултата на кампанията.

\paragraph{Правилото за собственост върху кадъра е правило за проектиране.} Формулировката от
раздел~\ref{sec:awaiter} е получена след като една и съща надпревара беше допусната три пъти
на три независими места. Следствието е, че интерфейсът трябва да прави грешката невъзможна,
и то важи независимо от това какво показва измерването.

\paragraph{Методиката е реализирана и самопроверена.} Инструментариумът съществува, изпълнява
правилата, които глава~V описва, и отказва да работи, когато условията му не са изпълнени.
Точно този отказ е и причината глава~VI да е празна, което е доказателство, че механизмът
работи, а не че липсва.

\paragraph{Решенията са използваеми в приложен код.} Глава~VII показва приложение, в което
handler-ите не съдържат условия по протокол, изгледите вземат данни от същите услуги като
API-то без излизане от процеса, а отложеното поле е добавка от един ред към обикновена
структура.

\subsection*{\bgen{Ограничения}{Limitations}}

Ограниченията се изброяват тук отново и накратко, защото заключение, което ги пропуска, ги
кани като въпроси.

Използваемостта е преценена по един източник, който е авторът на библиотеката. Независим
потребител не е привлечен, тоест глава~VII показва, че решенията могат да бъдат използвани,
а не че са удобни.

За HTTP/3 не съществува зрял генератор с постоянен темп на заявките в отворен цикъл.
Горните персентили за този протокол идват само от затворен цикъл и се означават като време
за обслужване, а не като време за отговор.

Под Windows липсва програмируем еквивалент на \code{tc}, а сегментирането на UDP се различава
съществено. Затова твърденията за производителност на HTTP/3 важат за Linux, а измерването
под Windows остава функционално.

Сравнението с насочване в ядрото не е проведено. Работата твърди преносимост на метода от
раздел~\ref{sec:cid} и измерва колко често препращането се налага, но не твърди превъзходство
над eBPF.

\subsection*{\bgen{Бъдеща работа}{Future work}}

Три посоки следват пряко от състоянието по-горе, и всяка от тях е дефинирана, а не пожелана.

\textbf{Провеждане на кампанията.} Изисква невиртуализирана машина под Linux с двойка мрежови
пространства. Планът, редът на изпълнение и правилата за допустимост са налични, тоест това е
изпълнение, а не проектиране.

\textbf{Рамо с насочване в ядрото.} След като делът на препратените пакети бъде измерен,
сравнението с програма на eBPF става смислено. До това измерване то би било сравнение без
известна величина от едната страна.

\textbf{Генератор с постоянен темп за HTTP/3.} Липсата му е ограничение на инструментариума в
целия отрасъл, не само тук, и премахването му би направило горните персентили за HTTP/3
сравними с тези за останалите два протокола.

\subsection*{\bgen{Обобщение}{Summary}}

Работата премества едно решение от конфигурацията на сокета към момента след приемането на
връзката, и показва, че това е достатъчно, за да отпадне зависимостта между броя на слушащите
дескриптори и броя на обслужваните протоколи. Постановката е доказана, реализацията работи по
трите протокола от един номер на порт, а методиката, която би определила цената, е налична и
самопроверена.

Остава цената да бъде измерена. Тя не е измерена и тук не се твърди каква е.

\TODO{числените изводи се дописват след провеждането на кампанията; дотогава заключението
съзнателно не съдържа стойност, за която няма измерване}
